\documentclass[a4paper]{report}

%%%%%%%%%%%%%%%%%%%%%%%%%%%%%%%%%%%%%%%%%%%%%%%%%%%%%%%%%%%%%%%%%%%%%%%%%%%%%%%%
% Provides page numbers and headings at the top of each page.
%%%%%%%%%%%%%%%%%%%%%%%%%%%%%%%%%%%%%%%%%%%%%%%%%%%%%%%%%%%%%%%%%%%%%%%%%%%%%%%%
\pagestyle{headings}

%%%%%%%%%%%%%%%%%%%%%%%%%%%%%%%%%%%%%%%%%%%%%%%%%%%%%%%%%%%%%%%%%%%%%%%%%%%%%%%%
% Packages
%%%%%%%%%%%%%%%%%%%%%%%%%%%%%%%%%%%%%%%%%%%%%%%%%%%%%%%%%%%%%%%%%%%%%%%%%%%%%%%%

% Warn about commands, classes and packages that are outdated and superseded.
% Orthodox checks for pitfalls that are not technically incorrect.
\usepackage[l2tabu,orthodox,abort]{nag}

% For font and input encoding.
\usepackage[utf8]{inputenc}
\usepackage[T1]{fontenc}

% Pretty much all of the ams maths packages.
\usepackage{amsmath,amsthm,amssymb,amsfonts}

% Allows inclusion of graphics easily and configurably.
\usepackage{graphicx}

% Provides commands to make subfigures.
\usepackage{subcaption}

% Typesets URLs sensibly --- with a TrueType font, hyperlinked in PDFs, and not
% breaking across lines.
\usepackage{url}

% Provides good access to colours.
\usepackage{color}
\usepackage{xcolor}

% Count the number of used registers (counter, dimen, skip, muskip, box, token,
% input, output, math families, languages, insertions). Also time the amount of
% time required for LaTeX compilation.
\usepackage[timer,left]{regstats}

\usepackage[noline,ruled]{algorithm2e}          % for algorithms
\usepackage{caption}                            % for subfloats
\usepackage[nodayofweek]{datetime}              % for \formatdate
\usepackage{float}                              % for setting float specifiers
\usepackage{../lib/gnuplot-lua-tikz/gnuplot-lua-tikz}                   % for plots produced by GNUplot
\usepackage{pdfpages}                           % for inserting compliance sheet as a PDF
\usepackage{pgfplots}                           % for matlab2tikz plots
\usepackage{subcaption}                         % for subfloats
\usepackage{tabularx}                           % for simple column stretching
\usepackage{tikz}                               % for graphics
\usetikzlibrary{positioning,shapes,shadows,arrows}
\usepackage{verbatim}                           % for comment environment

% My custom packages
\usepackage{escape}                             % for \escape
\usepackage{citationneeded}                     % for \citeNeeded
\usepackage{pieplots}                           % for \pieplots, \pieplot
\usepackage{pgf-pie}                            % for \pie

%%%%%%%%%%%%%%%%%%%%%%%%%%%%%%%%%%%%%%%%%%%%%%%%%%%%%%%%%%%%%%%%%%%%%%%%%%%%%%%%
% Source code formatting
%%%%%%%%%%%%%%%%%%%%%%%%%%%%%%%%%%%%%%%%%%%%%%%%%%%%%%%%%%%%%%%%%%%%%%%%%%%%%%%%
\usepackage{listings}

% Custom colours
\definecolor{darkGreen}{rgb}{0,0.6,0}
\definecolor{gray}{rgb}{0.5,0.5,0.5}
\definecolor{mauve}{rgb}{0.58,0,0.82}

% Default settings
\lstset{
    basicstyle=\ttfamily,           % font
                                    %
    title=\lstname,                 % show the filename of files included with \lstinputlisting
    captionpos=b,                   % sets the caption-position to bottom
    backgroundcolor=\color{white},  % choose the background color
                                    %
    numbers=left,                   % where to put the line-numbers
    numberstyle=\tiny\color{gray},  % the style that is used for the line-numbers
    stepnumber=1,                   % the step between two line-numbers
    numbersep=5pt,                  % how far the line-numbers are from the code
                                    %
    showspaces=false,               % show spaces adding particular underscores
    showstringspaces=false,         % underline spaces within strings
    showtabs=false,                 % show tabs within strings adding particular underscores
    tabsize=4,                      % sets default tabsize
    breaklines=true,                % sets automatic line breaking
    breakatwhitespace=false,        % sets if automatic breaks should only happen at whitespace
                                    %
    keywordstyle=\color{blue},      % keyword style
    commentstyle=\color{darkGreen}, % comment style
    stringstyle=\color{mauve},      % string literal style
    escapeinside={\%*}{*)},         % if you want to add a comment within your code
    morekeywords={*,...}            % if you want to add more keywords to the set
}
\lstloadlanguages{C,Matlab}

%%%%%%%%%%%%%%%%%%%%%%%%%%%%%%%%%%%%%%%%%%%%%%%%%%%%%%%%%%%%%%%%%%%%%%%%%%%%%%%%
% Gantt charts
%%%%%%%%%%%%%%%%%%%%%%%%%%%%%%%%%%%%%%%%%%%%%%%%%%%%%%%%%%%%%%%%%%%%%%%%%%%%%%%%
\usepackage{tikz}
\usetikzlibrary{calendar}
\usepackage{pgfgantt}

%%%%%%%%%%%%%%%%%%%%%%%%%%%%%%%%%%%%%%%%%%%%%%%%%%%%%%%%%%%%%%%%%%%%%%%%%%%%%%%%
% Variables (to save retyping and reformatting these throughout the document)
%%%%%%%%%%%%%%%%%%%%%%%%%%%%%%%%%%%%%%%%%%%%%%%%%%%%%%%%%%%%%%%%%%%%%%%%%%%%%%%%
%%%%%%%%%%%%%%%%%%%%%%%%%%%%%%%%%%%%%%%%%%%%%%%%%%%%%%%%%%%%%%%%%%%%%%%%%%%%%%%%
% Variables (to save retyping and reformatting these throughout the document).
%%%%%%%%%%%%%%%%%%%%%%%%%%%%%%%%%%%%%%%%%%%%%%%%%%%%%%%%%%%%%%%%%%%%%%%%%%%%%%%%
\usepackage{definevariable} % for \defineVariable
\usepackage{persons} % for \definePerson, \getPerson, \formatPerson

\defineVariable\thesis{}{Thesis}
\defineVariable\thesisTitle{\textsc}{Implementation of Random Projection Algorithms Using an FPGA}
\defineVariable\thesisDate{}{2012}
\defineVariable\studentName{\formatPerson}{\getPerson*{Spence}}
\defineVariable\supervisorName{\formatPerson}{\getPerson*{Leong}}

\defineVariable\degreeName{}{Bachelor of Engineering (Computer)}
\defineVariable\schoolName{}{School of Electrical Engineering}
\defineVariable\studentID{}{308216350}
\defineVariable\universityName{}{The University of Sydney}
\defineVariable\uosAssessment{}{\thesis}
\defineVariable\uosCode{}{ELEC4712/ELEC4713}

\urldef{\gitRepoHTTP}\url{https://github.com/joshuaspence/ThesisCode.git}
\urldef{\gitRepoSSH}\url{git@github.com:joshuaspence/ThesisCode.git}

% People
\definePerson{Chawla}{Sanjay Chawla}
\definePerson{Frechtling}{Michael Frechtling}
\definePerson{Heinrichs}{Amelia Heinrichs}
\definePerson{Khoa}{Nguyen Lu Dang Khoa}
\definePerson{Leong}{Philip Leong}
\definePerson{Spence}{Joshua Spence}

% PDF metadata
\newcommand*\pdfmetadata{%
    pdfauthor={Joshua Spence},
    pdftitle={Implementation of Random Projection Algorithms Using an FPGA},
    pdfsubject={},
    pdfkeywords={commute time,anomaly detection,FPGA},
    pdfdisplaydoctitle={true}
}

% Copyrights and trademarks
\newcommand\IntelXeon{Intel\textsuperscript{\textregistered} Xeon\textsuperscript{\textregistered}}
\newcommand\IntelCoreiSeven{Intel\textsuperscript{\textregistered} Core\textsuperscript{\texttrademark} i7}
\newcommand\AutoESL{\software{Xilinx\textsuperscript{\textregistered} AutoESL}}
\newcommand\ModelSim{\software{Altera\textsuperscript{\textregistered} ModelSim}}
\newcommand\Xilinx{Xilinx\textsuperscript{\textregistered}}
\newcommand\ARM{ARM\textsuperscript{\textregistered}}
\newcommand\Cortex{Cortex\textsuperscript{\texttrademark}}

%%%%%%%%%%%%%%%%%%%%%%%%%%%%%%%%%%%%%%%%%%%%%%%%%%%%%%%%%%%%%%%%%%%%%%%%%%%%%%%%
% Title page options
%%%%%%%%%%%%%%%%%%%%%%%%%%%%%%%%%%%%%%%%%%%%%%%%%%%%%%%%%%%%%%%%%%%%%%%%%%%%%%%%
\ifcsname beamer@author\endcsname
    \title[Random Projection Algorithms Using an FPGA]{Implementation of Random Projection Algorithms Using an FPGA}
    \author[Joshua Spence]{Joshua Spence\\{\tiny{Supervisor: Philip Leong}}}
    \institute[The University of Sydney]{
        School of Electrical Engineering\\
        The University of Sydney
    }
    \titlegraphic{\includegraphics[width=0.25\textwidth]{sydney_uni_coat_of_arms}}
\else
    \title{Implementation of Random Projection Algorithms Using an FPGA}
    \author{Joshua Spence}
\fi

\date{2012}

%%%%%%%%%%%%%%%%%%%%%%%%%%%%%%%%%%%%%%%%%%%%%%%%%%%%%%%%%%%%%%%%%%%%%%%%%%%%%%%%
% Custom commands
%%%%%%%%%%%%%%%%%%%%%%%%%%%%%%%%%%%%%%%%%%%%%%%%%%%%%%%%%%%%%%%%%%%%%%%%%%%%%%%%
\usepackage{listings} % for \lstinline
\usepackage{xparse} % for \NewDocumentCommand

% To allow us to easily format a table row.
% See http://www.latex-community.org/forum/viewtopic.php?f=5&t=3293
\newcolumntype{+}{>{\global\let\currentrowstyle\relax}}
\newcolumntype{^}{>{\currentrowstyle}}
\newcommand\rowstyle[1]{\gdef\currentrowstyle{#1}%
	#1\ignorespaces%
}

% \tableHeader{...}
\newcommand\tableHeader[1]{%
	\hline \rowstyle{\bfseries} #1 \\\hline%
}

% \command{CMD}
\newcommand\command[1]{\lstinline`#1`}

% \dataset{DATASET}
\newcommand\dataset[1]{`\escape{#1}'}

% \hostname{HOSTNAME}
\newcommand\hostname[1]{\emph{#1}}

% \programmingLanguage{LANG}
\newcommand\programmingLanguage[1]{\lstinline`#1`}

% \software{NAME}[CMD]
\NewDocumentCommand\software{om}{``#2'' \IfNoValueTF{#1}{}{(\command{#1})}}

% Simple command I defined to allow me to mark TODO items in red
\newcommand{\todo}[1] {\textbf{\textcolor{red}{#1}}}

%%%%%%%%%%%%%%%%%%%%%%%%%%%%%%%%%%%%%%%%%%%%%%%%%%%%%%%%%%%%%%%%%%%%%%%%%%%%%%%%
% \begin{specifications}
% \specificationGroup{GROUP_NAME}
% \specification{NAME}{VALUE}
% \endspecificationGroup
% \end{specifications}
%%%%%%%%%%%%%%%%%%%%%%%%%%%%%%%%%%%%%%%%%%%%%%%%%%%%%%%%%%%%%%%%%%%%%%%%%%%%%%%%
\usepackage{xparse} % for \NewDocumentEnvironment

\makeatletter
\NewDocumentEnvironment{specifications}{}{%
    \def\@hline{}%
    %
    \newenvironment{specificationGroup}[1]{%
        \@hline \underline{##1} &%
        \def\@hline{\\}%
    }{%
        \def\@hline{\\[1em]}%
    }%
    \newcommand\specification[2]{%
        \@hline ##1 & ##2%
        \def\@hline{\\}%
    }%
    %
    \begin{tabular}{|+>{\bfseries}l|^p{8cm}|}%
        \tableHeader{Item & Specification}%
}{%
        \\\hline%
    \end{tabular}%
}
\makeatother

%%%%%%%%%%%%%%%%%%%%%%%%%%%%%%%%%%%%%%%%%%%%%%%%%%%%%%%%%%%%%%%%%%%%%%%%%%%%%%%%
% \begin{datasetDescription}{DATASET}
% \datasetProperty{NAME}{VALUE}
% \end{datasetDescription}
%%%%%%%%%%%%%%%%%%%%%%%%%%%%%%%%%%%%%%%%%%%%%%%%%%%%%%%%%%%%%%%%%%%%%%%%%%%%%%%%
\usepackage{escape} % for \escape
\usepackage{graphicx} % for \includegraphics
\usepackage{xparse} % for \NewDocumentEnvironment

\makeatletter
\NewDocumentEnvironment{datasetDescription}{m}{%
    \def\@hline{}%
    %
    \newcommand\datasetProperty[2]{%
        \@hline ##1 & ##2%
        \def\@hline{\\}%
    }%
    %
    \def\tablewidth{8cm}%
    \begin{longtable}{|+>{\bfseries}l|^p{\tablewidth}|}%
        \tableHeader{Property & Value}%
        \datasetProperty{Identifier}{\escape{#1}}%
}{%
        \\\hline%
    \end{longtable}%
    %
    % NOTE: The "plots/datasets/#1" file can be quite large. You may need to
    % inrease the amount of memory that the LaTeX compiler has available for
    % use. Basically just pdate the value of "main_memory" in the 
    % /etc/texmf/texmf.cnf file.
    %
    % See: http://tex.stackexchange.com/questions/27529/increase-tex-capacity-as-non-root
    \begin{figure}[H]%
        \centering%
        \includegraphics[width=0.5\textwidth]{pca/#1.png}%
        \caption{A PCA plot for the `\escape{#1}' data set}%
        \label{dataSets:#1:pca}%
    \end{figure}%
}
\makeatother

%%%%%%%%%%%%%%%%%%%%%%%%%%%%%%%%%%%%%%%%%%%%%%%%%%%%%%%%%%%%%%%%%%%%%%%%%%%%%%%%
% Create a blank page
%%%%%%%%%%%%%%%%%%%%%%%%%%%%%%%%%%%%%%%%%%%%%%%%%%%%%%%%%%%%%%%%%%%%%%%%%%%%%%%%
\newcommand{\blankpage}{%
    \newpage%
    \thispagestyle{empty}%
    \mbox{}%
    \newpage%
}

% Makes references hyperlinks in PDF output. Also allows automatic references
% using \autoref.
%
% Note that this package must be almost last in the preamble. In particular, it
% must appear *after* "float" and *before* "algorithm".
\usepackage[breaklinks=true,hidelinks,\pdfmetadata]{hyperref}

%%%%%%%%%%%%%%%%%%%%%%%%%%%%%%%%%%%%%%%%%%%%%%%%%%%%%%%%%%%%%%%%%%%%%%%%%%%%%%%%
% Graphics directory
%%%%%%%%%%%%%%%%%%%%%%%%%%%%%%%%%%%%%%%%%%%%%%%%%%%%%%%%%%%%%%%%%%%%%%%%%%%%%%%%
\graphicspath{{../img/}}

%%%%%%%%%%%%%%%%%%%%%%%%%%%%%%%%%%%%%%%%%%%%%%%%%%%%%%%%%%%%%%%%%%%%%%%%%%%%%%%%
% Title page options
%%%%%%%%%%%%%%%%%%%%%%%%%%%%%%%%%%%%%%%%%%%%%%%%%%%%%%%%%%%%%%%%%%%%%%%%%%%%%%%%
\title{\thesisTitle*{}}
\author{\studentName*{}}
\date{\thesisDate*{}}

%%%%%%%%%%%%%%%%%%%%%%%%%%%%%%%%%%%%%%%%%%%%%%%%%%%%%%%%%%%%%%%%%%%%%%%%%%%%%%%%
% Glossary
%%%%%%%%%%%%%%%%%%%%%%%%%%%%%%%%%%%%%%%%%%%%%%%%%%%%%%%%%%%%%%%%%%%%%%%%%%%%%%%%
\usepackage[toc,xindy]{glossaries} % note that hyperref must be loaded first
\makeglossaries
\newacronym{ASIC}{ASIC}{Application-Specific Integrated Circuit}
\newacronym{EPROM}{EPROM}{Erasable Programmable Read Only Memory}
\newacronym{FPGA}{FPGA}{Field-Programmable Gate Arrays}
\newacronym{HDL}{HDL}{Hardware Description Language}
\newacronym{LUT}{LUT}{Lookup Table}
\newacronym{PCA}{PCA}{Principal Component Analysis}
\newacronym{RTL}{RTL}{Register Transfer Level}
\newacronym{SRAM}{SRAM}{Static Random Access Memory}
\newacronym{VHDL}{VHDL}{VHSIC Hardware Description Language}


%%%%%%%%%%%%%%%%%%%%%%%%%%%%%%%%%%%%%%%%%%%%%%%%%%%%%%%%%%%%%%%%%%%%%%%%%%%%%%%%
% Bibliography
%%%%%%%%%%%%%%%%%%%%%%%%%%%%%%%%%%%%%%%%%%%%%%%%%%%%%%%%%%%%%%%%%%%%%%%%%%%%%%%%
\usepackage[language=australian]{biblatex} % use biblatex instead of bibtex
\bibliography{thesis}

%%%%%%%%%%%%%%%%%%%%%%%%%%%%%%%%%%%%%%%%%%%%%%%%%%%%%%%%%%%%%%%%%%%%%%%%%%%%%%%%
\begin{document}

\pagestyle{empty} % suppress page numbers

%%%%%%%%%%%%%%%%%%%%%%%%%%%%%%%%%%%%%%%%%%%%%%%%%
% Title Page
%%%%%%%%%%%%%%%%%%%%%%%%%%%%%%%%%%%%%%%%%%%%%%%%%
\pdfbookmark[0]{Preamble}{title}
\pdfbookmark[1]{Title page}{title}
\begin{titlepage}
\clearpage
\label{title}
\centering

% Upper part of the page
\includegraphics[width=0.25\textwidth]{sydney_uni_coat_of_arms}\\[1cm]
\textsc{\LARGE \universityName*{}}\\[1.5cm]
\textsc{\Large Final year thesis}\\[0.5cm]

% Title
\rule{\linewidth}{0.5mm}\\[0.4cm] % horiizontal line
{\huge \bfseries \thesisTitle*{}}\\[0.4cm]
\rule{\linewidth}{0.5mm}\\[1.5cm] % horizontal line

% Author and supervisor
\begin{minipage}{0.4\textwidth}
    \begin{flushleft}
        \large
        \emph{Author:}\\
        \studentName*{}
    \end{flushleft}
\end{minipage}
\begin{minipage}{0.4\textwidth}
    \begin{flushright}
        \large
        \emph{Supervisor:}\\
        \supervisorName*{}
    \end{flushright}
\end{minipage}

% Strech vertical space so that it fills all empty space
\vfill

% Thesis description
{\large A thesis submitted in fulfilment of the requirements for the\\
degree of Bachelor of Engineering (Computer) in the\\
School of Electrical Engineering at\\
The University of Sydney}

% Strech vertical space so that it fills all empty space
\vfill

% Bottom of the page
{\large \thesisDate*{}}

\end{titlepage}

%%%%%%%%%%%%%%%%%%%%%%%%%%%%%%%%%%%%%%%%%%%%%%%%%
% Compliance Sheet
%%%%%%%%%%%%%%%%%%%%%%%%%%%%%%%%%%%%%%%%%%%%%%%%%
\clearpage
\pdfbookmark[1]{Compliance Sheet}{compliance}
\label{compliance}
\includepdf[
        scale=0.8,
        pagecommand={
            \begin{tikzpicture}[
                    remember picture,
                    overlay,
                    shift=(current page.center)
                    ]
                \tikzstyle{every node}=[font=\large]
                \coordinate (name-pos)          at (-5.4, -1.6);
                \coordinate (sid-pos)           at (-5.4, -3.0);
                \coordinate (date-pos)          at (-5.4, -5.8);
                \coordinate (uos-pos)           at (-7.4, -7.7);
                \coordinate (assessment-pos)    at (-7.4, -9.8);
                \node[anchor=west] (name)       at (name-pos)       {\studentName*{}};
                \node[anchor=west] (sid)        at (sid-pos)        {\studentID*{}};
                \node[anchor=west] (date)       at (date-pos)       {\today{}};
                \node[anchor=west] (uos)        at (uos-pos)        {\uosCode*{}};
                \node[anchor=west] (assessment) at (assessment-pos) {\uosAssessment*{}};
            \end{tikzpicture}
        }
    ]{../forms/compliance_sheet}

\begin{comment}
%%%%%%%%%%%%%%%%%%%%%%%%%%%%%%%%%%%%%%%%%%%%%%%%%
% Statement of Achievement
%%%%%%%%%%%%%%%%%%%%%%%%%%%%%%%%%%%%%%%%%%%%%%%%%
\clearpage
\pdfbookmark[1]{Statement of Achievement}{statement-of-achievement}
\chapter*{Statement of Achievement}
\label{statement-of-achievement}
% This is located in {chapter01/contributions}
\end{comment}

%%%%%%%%%%%%%%%%%%%%%%%%%%%%%%%%%%%%%%%%%%%%%%%%%
% Acknowledgements
%%%%%%%%%%%%%%%%%%%%%%%%%%%%%%%%%%%%%%%%%%%%%%%%%
\clearpage
\pdfbookmark[1]{Acknowledgements}{acknowledgements}
\chapter*{Acknowledgements}
\label{acknowledgements}
This project could not have been completed without the guidance of my project
supervisor, \supervisorName{}. He always made himself available to assist me
with the understanding of many concepts, which were both new and unknown. His
expertise in the field was a great asset and greatly contributed to my own
understanding of the ideas explored in this thesis. His support and guidance was
invaluable, and were instrumental to the success of my project.

Much of the works completed during this project could not have been performed
without the prior research of \getPerson{Chawla} and \getPerson{Khoa}. Both were
generous enough to lend their own support and expertise in order to provide an
understanding of the implemented anomaly detection algorithm.

Additionally, I must thank \getPerson{Frechtling} for his assistance and advice
throughout the project. He assisted me with software packages including
\software{Xilinx\textsuperscript{\textregistered} AutoESL} and
\software{Altera\textsuperscript{\textregistered} ModelSim} (along with many
others), as well as offering valuable advice regarding my future professional
career.

Last but not least, none of this would have been possible without the support
and care of my family. My immediate family, to whom I deeply appreciate, always
has been a continuous source of support and encouragement throughout my
university degree. And I would like to express my gratitude to my girlfriend,
\getPerson{Heinrichs}, for her love, patience and support throughout the final
year of my university studies.

%%%%%%%%%%%%%%%%%%%%%%%%%%%%%%%%%%%%%%%%%%%%%%%%%
% Publications
%%%%%%%%%%%%%%%%%%%%%%%%%%%%%%%%%%%%%%%%%%%%%%%%%
\clearpage
\pdfbookmark[1]{Publications}{publications}
\chapter*{Publications}
\label{publications}
This Thesis is based on the following publications:
\begin{enumerate}
    \item \fullcite{Vries:2010}
    \item \fullcite{Khoa:2012}
    \item \fullcite{Bay:2003}
\end{enumerate}

%%%%%%%%%%%%%%%%%%%%%%%%%%%%%%%%%%%%%%%%%%%%%%%%%
% Abstract
%%%%%%%%%%%%%%%%%%%%%%%%%%%%%%%%%%%%%%%%%%%%%%%%%
\clearpage
\pdfbookmark[1]{Abstract}{abstract}
\begin{abstract}
\label{abstract}
The prediction of the stock market has become an issue of great interest in the
areas of finance, mathematics and engineering; due mainly to the great potential
financial gain. Researchers have devised various algorithms and differing
approaches to the problem of stock market analysis, with varying degrees of
success.

A major outstanding issue for stock market analysis is the effective and
efficient detection of local anomalies in the input data sets, which are
inherently highly multidimensional. Many na\"{\i}ve algorithms are highly
inefficient and others fail to adequately detect local anomalies altogether. It
had become a time-vs-correctness trade-off in which no acceptable compromise
could be reached.

However, researchers are starting to explore the relatively new concept of
applying ``random projections'' to the highly multidimensional data sets.
Research has suggested that by applying these random projections, they are able
to significantly reduce the dimensionality (and consequently the computational
complexity) of the data sets, whilst sufficiently retaining the inherent
properties of that data set --- at least so much so as anomaly detection is
concerned.

Anomaly detection is important because it allows otherwise-accurate machine
learning algorithms such as neural networks to more accurately model and predict
the stock exchange data by ignoring anomalous data, which likely doesn't effect
the state of the model to any significant degree.

\getPerson{Chawla} from \universityName{} has in recent years conducted and
supervised new and exciting research oriented around random projections. In
particular, \getPerson{Khoa}, under the supervision of \getPerson{Chawla}
evaluated the use of traditional distance metrics, such as Euclidean distance
and Mahalanobis distance, in the application of local anomaly detection.
\getPerson{Khoa} proposed the use of the `commute time' metric, derived from
random walks on graphs, in anomaly detection.

The source code that was used throughout this project can be freely downloaded
from \gitRepoHTTP{} (over HTTP) or \gitRepoSSH{} (over SSH).
\end{abstract}

\pagenumbering{roman}

%%%%%%%%%%%%%%%%%%%%%%%%%%%%%%%%%%%%%%%%%%%%%%%%%
% Table of Contents
%%%%%%%%%%%%%%%%%%%%%%%%%%%%%%%%%%%%%%%%%%%%%%%%%
\clearpage
\pdfbookmark[0]{Contents}{toc}
\renewcommand\contentsname{Table of Contents}
\pdfbookmark[1]{\contentsname}{toc}
\tableofcontents

%%%%%%%%%%%%%%%%%%%%%%%%%%%%%%%%%%%%%%%%%%%%%%%%%
% List of Figures
%%%%%%%%%%%%%%%%%%%%%%%%%%%%%%%%%%%%%%%%%%%%%%%%%
\clearpage
\pdfbookmark[1]{\listfigurename}{lof}
\listoffigures

%%%%%%%%%%%%%%%%%%%%%%%%%%%%%%%%%%%%%%%%%%%%%%%%%
% List of Algorithms
%%%%%%%%%%%%%%%%%%%%%%%%%%%%%%%%%%%%%%%%%%%%%%%%%
\clearpage
\pdfbookmark[1]{\listalgorithmcfname}{loa}
\listofalgorithms

%%%%%%%%%%%%%%%%%%%%%%%%%%%%%%%%%%%%%%%%%%%%%%%%%
% List of Tables
%%%%%%%%%%%%%%%%%%%%%%%%%%%%%%%%%%%%%%%%%%%%%%%%%
\clearpage
\pdfbookmark[1]{\listtablename}{lot}
\listoftables

\blankpage
\pagenumbering{arabic}

%%%%%%%%%%%%%%%%%%%%%%%%%%%%%%%%%%%%%%%%%%%%%%%%%
% CHAPTER 01: Introduction
%%%%%%%%%%%%%%%%%%%%%%%%%%%%%%%%%%%%%%%%%%%%%%%%%
\chapter{Introduction}
\label{introduction}
\chapter{Introduction}
\label{introduction}

%%%%%%%%%%%%%%%%%%%%%%%%%%%%%%%%%%%%%%%%%%%%%%%%%
% Motivation
%%%%%%%%%%%%%%%%%%%%%%%%%%%%%%%%%%%%%%%%%%%%%%%%%
\section{Motivation}
\label{motivation}
Anomaly detection is an important technique which can be applied to a wide range
of applications. There are many techniques to detect and measure anomalies, each
usually most applicable to some specific problem domain \citeNeeded{}. A general
algorithm for anomaly detection has proved difficult to design \citeNeeded{},
due to various challenges associated with defining and measuring anomalies. Most
existing approaches are based on statistical or geometrical measures involving
distance. Whilst many algorithms work well for some subset of input data sets,
it has been difficult to discover an algorithm which performs both correctly and
efficiently on all possible data sets.

Previous research has found merit in applying randomization techniques to highly
multivariate data sets in order to reduce the dimensionality of these data sets
whilst maintaining their fundamental and statistical properties. Such reduction
of the dimensionality of data, assuming it can be performed efficiently, allows
previously-unscalable anomaly detection algorithms to be practically applied to
a wider range of data and applications.

Anomaly detection is an important and contemporary problem in the field of
computer science, and is of particular interesting to stock market analysis,
network intrusion detection and image comparison.

%%%%%%%%%%%%%%%%%%%%%%%%%%%%%%%%%%%%%%%%%%%%%%%%%
% Contributions of this thesis
%%%%%%%%%%%%%%%%%%%%%%%%%%%%%%%%%%%%%%%%%%%%%%%%%
\section{Contributions of this thesis}
\label{contributions}
One key difficulty in anomaly detection is the efficient scaling of a general
algorithm to apply to highly multivariate data \citeNeeded{}. In this thesis, I
explore the use of randomization techniques (such as random projections and
commute time) in anomaly detection algorithms. These techniques provide
encouraging results with regards to the run-time complexity of an algorithm.

Furthermore, in this thesis I attempt to make observations and an analysis of
the run-time performance of anomaly detection algorithms using randomization
techniques, so as to identify steps in the algorithms that bottleneck the
algorithm's performance. Through this identification it would be possible to
improve the \emph{actual} run-time performance of the algorithm by utilising
the advantages of reconfigurable computing.


%%%%%%%%%%%%%%%%%%%%%%%%%%%%%%%%%%%%%%%%%%%%%%%%%
% Organization
%%%%%%%%%%%%%%%%%%%%%%%%%%%%%%%%%%%%%%%%%%%%%%%%%
\section{Organization}
\label{organization}
The rest of this \thesis{} is organized as follows. In \autoref{background}, I
provide a background to various anomaly detection techniques, as well as
randomization techniques. In order to provide the reader with an understanding
of the background topics, a brief background is given to various topics in
linear algebra, vector calculus and graph theory. However, the provided
background material is not intended to be comprehensive, and it is expected that
the reader is familiar with these fields of mathematics. In
\autoref{reconfigurableComputing}, I provide an overview of reconfigurable
computing, including an explanation of \glspl{FPGA} and an overview of
\gls{HLS}.

In \autoref{software}, I profile the execution of an anomaly detection algorithm
and explore possible improvements to the algorithm, in particular by outsourcing
various stages of the algorithm to an FPGA device. I proceed to make iterative
refinements to the algorithm from its base implementation towards a hardware
design.

In \autoref{hardware}, I propose several hardware implementations and discuss
design considerations, as well as the design process itself, that must be
considered for these designs. I also detail the results of some initial testing
investigating the performance expectations from differing hardware designs.

I conclude in \autoref{conclusions} by reflecting upon the results obtained
through this research, and making suggestions for further research in this
topic.

 
%%%%%%%%%%%%%%%%%%%%%%%%%%%%%%%%%%%%%%%%%%%%%%%%%
% Schedule
%%%%%%%%%%%%%%%%%%%%%%%%%%%%%%%%%%%%%%%%%%%%%%%%%
\section{Schedule}
\label{schedule}
A Gantt chart showing the anticipated schedule for the project is shown in 
\autoref{fig:ganttChart}. This will be updated as the project progresses.

% Gantt chart
% See http://www.martin-kumm.de/tex_gantt_package.php
\begin{figure}[h]
	\centering
	\scalebox{0.5}{
	\begin{gantt}[
			xunitlength=0.5cm,
			fontsize=\small,
			titlefontsize=\small,
			drawledgerline=true
			]{18}{36}
		% Months
		\begin{ganttitle}
			\titleelement{Mar}{4}
			\titleelement{Apr}{4}
			\titleelement{May}{4}
			\titleelement{Jun}{4}
			\titleelement{Jul}{4}
			\titleelement{Aug}{4}
			\titleelement{Sep}{4}
			\titleelement{Oct}{4}
			\titleelement{Nov}{4}
		\end{ganttitle}
		% Split months into quarters
		\begin{ganttitle}
			\numtitle{1}{1}{4}{1}
			\numtitle{1}{1}{4}{1}
			\numtitle{1}{1}{4}{1}
			\numtitle{1}{1}{4}{1}
			\numtitle{1}{1}{4}{1}
			\numtitle{1}{1}{4}{1}
			\numtitle{1}{1}{4}{1}
			\numtitle{1}{1}{4}{1}
			\numtitle{1}{1}{4}{1}
		\end{ganttitle}

		\input{ganttdata}
	\end{gantt}
	}
	\caption{Schedule for thesis work}
	\label{fig:ganttChart}
\end{figure}



%%%%%%%%%%%%%%%%%%%%%%%%%%%%%%%%%%%%%%%%%%%%%%%%%
% CHAPTER 02: Background
%%%%%%%%%%%%%%%%%%%%%%%%%%%%%%%%%%%%%%%%%%%%%%%%%
\chapter{Background}
\label{background}
\chapter{Introduction}
\label{introduction}

%%%%%%%%%%%%%%%%%%%%%%%%%%%%%%%%%%%%%%%%%%%%%%%%%
% Motivation
%%%%%%%%%%%%%%%%%%%%%%%%%%%%%%%%%%%%%%%%%%%%%%%%%
\section{Motivation}
\label{motivation}
Anomaly detection is an important technique which can be applied to a wide range
of applications. There are many techniques to detect and measure anomalies, each
usually most applicable to some specific problem domain \citeNeeded{}. A general
algorithm for anomaly detection has proved difficult to design \citeNeeded{},
due to various challenges associated with defining and measuring anomalies. Most
existing approaches are based on statistical or geometrical measures involving
distance. Whilst many algorithms work well for some subset of input data sets,
it has been difficult to discover an algorithm which performs both correctly and
efficiently on all possible data sets.

Previous research has found merit in applying randomization techniques to highly
multivariate data sets in order to reduce the dimensionality of these data sets
whilst maintaining their fundamental and statistical properties. Such reduction
of the dimensionality of data, assuming it can be performed efficiently, allows
previously-unscalable anomaly detection algorithms to be practically applied to
a wider range of data and applications.

Anomaly detection is an important and contemporary problem in the field of
computer science, and is of particular interesting to stock market analysis,
network intrusion detection and image comparison.

%%%%%%%%%%%%%%%%%%%%%%%%%%%%%%%%%%%%%%%%%%%%%%%%%
% Contributions of this thesis
%%%%%%%%%%%%%%%%%%%%%%%%%%%%%%%%%%%%%%%%%%%%%%%%%
\section{Contributions of this thesis}
\label{contributions}
One key difficulty in anomaly detection is the efficient scaling of a general
algorithm to apply to highly multivariate data \citeNeeded{}. In this thesis, I
explore the use of randomization techniques (such as random projections and
commute time) in anomaly detection algorithms. These techniques provide
encouraging results with regards to the run-time complexity of an algorithm.

Furthermore, in this thesis I attempt to make observations and an analysis of
the run-time performance of anomaly detection algorithms using randomization
techniques, so as to identify steps in the algorithms that bottleneck the
algorithm's performance. Through this identification it would be possible to
improve the \emph{actual} run-time performance of the algorithm by utilising
the advantages of reconfigurable computing.


%%%%%%%%%%%%%%%%%%%%%%%%%%%%%%%%%%%%%%%%%%%%%%%%%
% Organization
%%%%%%%%%%%%%%%%%%%%%%%%%%%%%%%%%%%%%%%%%%%%%%%%%
\section{Organization}
\label{organization}
The rest of this \thesis{} is organized as follows. In \autoref{background}, I
provide a background to various anomaly detection techniques, as well as
randomization techniques. In order to provide the reader with an understanding
of the background topics, a brief background is given to various topics in
linear algebra, vector calculus and graph theory. However, the provided
background material is not intended to be comprehensive, and it is expected that
the reader is familiar with these fields of mathematics. In
\autoref{reconfigurableComputing}, I provide an overview of reconfigurable
computing, including an explanation of \glspl{FPGA} and an overview of
\gls{HLS}.

In \autoref{software}, I profile the execution of an anomaly detection algorithm
and explore possible improvements to the algorithm, in particular by outsourcing
various stages of the algorithm to an FPGA device. I proceed to make iterative
refinements to the algorithm from its base implementation towards a hardware
design.

In \autoref{hardware}, I propose several hardware implementations and discuss
design considerations, as well as the design process itself, that must be
considered for these designs. I also detail the results of some initial testing
investigating the performance expectations from differing hardware designs.

I conclude in \autoref{conclusions} by reflecting upon the results obtained
through this research, and making suggestions for further research in this
topic.

 
%%%%%%%%%%%%%%%%%%%%%%%%%%%%%%%%%%%%%%%%%%%%%%%%%
% Schedule
%%%%%%%%%%%%%%%%%%%%%%%%%%%%%%%%%%%%%%%%%%%%%%%%%
\section{Schedule}
\label{schedule}
A Gantt chart showing the anticipated schedule for the project is shown in 
\autoref{fig:ganttChart}. This will be updated as the project progresses.

% Gantt chart
% See http://www.martin-kumm.de/tex_gantt_package.php
\begin{figure}[h]
	\centering
	\scalebox{0.5}{
	\begin{gantt}[
			xunitlength=0.5cm,
			fontsize=\small,
			titlefontsize=\small,
			drawledgerline=true
			]{18}{36}
		% Months
		\begin{ganttitle}
			\titleelement{Mar}{4}
			\titleelement{Apr}{4}
			\titleelement{May}{4}
			\titleelement{Jun}{4}
			\titleelement{Jul}{4}
			\titleelement{Aug}{4}
			\titleelement{Sep}{4}
			\titleelement{Oct}{4}
			\titleelement{Nov}{4}
		\end{ganttitle}
		% Split months into quarters
		\begin{ganttitle}
			\numtitle{1}{1}{4}{1}
			\numtitle{1}{1}{4}{1}
			\numtitle{1}{1}{4}{1}
			\numtitle{1}{1}{4}{1}
			\numtitle{1}{1}{4}{1}
			\numtitle{1}{1}{4}{1}
			\numtitle{1}{1}{4}{1}
			\numtitle{1}{1}{4}{1}
			\numtitle{1}{1}{4}{1}
		\end{ganttitle}

		\input{ganttdata}
	\end{gantt}
	}
	\caption{Schedule for thesis work}
	\label{fig:ganttChart}
\end{figure}



%%%%%%%%%%%%%%%%%%%%%%%%%%%%%%%%%%%%%%%%%%%%%%%%%
% CHAPTER 03: Reconfigurable Computing
%%%%%%%%%%%%%%%%%%%%%%%%%%%%%%%%%%%%%%%%%%%%%%%%%
\chapter{Reconfigurable Computing}
\label{reconfigurableComputing}
\chapter{Introduction}
\label{introduction}

%%%%%%%%%%%%%%%%%%%%%%%%%%%%%%%%%%%%%%%%%%%%%%%%%
% Motivation
%%%%%%%%%%%%%%%%%%%%%%%%%%%%%%%%%%%%%%%%%%%%%%%%%
\section{Motivation}
\label{motivation}
Anomaly detection is an important technique which can be applied to a wide range
of applications. There are many techniques to detect and measure anomalies, each
usually most applicable to some specific problem domain \citeNeeded{}. A general
algorithm for anomaly detection has proved difficult to design \citeNeeded{},
due to various challenges associated with defining and measuring anomalies. Most
existing approaches are based on statistical or geometrical measures involving
distance. Whilst many algorithms work well for some subset of input data sets,
it has been difficult to discover an algorithm which performs both correctly and
efficiently on all possible data sets.

Previous research has found merit in applying randomization techniques to highly
multivariate data sets in order to reduce the dimensionality of these data sets
whilst maintaining their fundamental and statistical properties. Such reduction
of the dimensionality of data, assuming it can be performed efficiently, allows
previously-unscalable anomaly detection algorithms to be practically applied to
a wider range of data and applications.

Anomaly detection is an important and contemporary problem in the field of
computer science, and is of particular interesting to stock market analysis,
network intrusion detection and image comparison.

%%%%%%%%%%%%%%%%%%%%%%%%%%%%%%%%%%%%%%%%%%%%%%%%%
% Contributions of this thesis
%%%%%%%%%%%%%%%%%%%%%%%%%%%%%%%%%%%%%%%%%%%%%%%%%
\section{Contributions of this thesis}
\label{contributions}
One key difficulty in anomaly detection is the efficient scaling of a general
algorithm to apply to highly multivariate data \citeNeeded{}. In this thesis, I
explore the use of randomization techniques (such as random projections and
commute time) in anomaly detection algorithms. These techniques provide
encouraging results with regards to the run-time complexity of an algorithm.

Furthermore, in this thesis I attempt to make observations and an analysis of
the run-time performance of anomaly detection algorithms using randomization
techniques, so as to identify steps in the algorithms that bottleneck the
algorithm's performance. Through this identification it would be possible to
improve the \emph{actual} run-time performance of the algorithm by utilising
the advantages of reconfigurable computing.


%%%%%%%%%%%%%%%%%%%%%%%%%%%%%%%%%%%%%%%%%%%%%%%%%
% Organization
%%%%%%%%%%%%%%%%%%%%%%%%%%%%%%%%%%%%%%%%%%%%%%%%%
\section{Organization}
\label{organization}
The rest of this \thesis{} is organized as follows. In \autoref{background}, I
provide a background to various anomaly detection techniques, as well as
randomization techniques. In order to provide the reader with an understanding
of the background topics, a brief background is given to various topics in
linear algebra, vector calculus and graph theory. However, the provided
background material is not intended to be comprehensive, and it is expected that
the reader is familiar with these fields of mathematics. In
\autoref{reconfigurableComputing}, I provide an overview of reconfigurable
computing, including an explanation of \glspl{FPGA} and an overview of
\gls{HLS}.

In \autoref{software}, I profile the execution of an anomaly detection algorithm
and explore possible improvements to the algorithm, in particular by outsourcing
various stages of the algorithm to an FPGA device. I proceed to make iterative
refinements to the algorithm from its base implementation towards a hardware
design.

In \autoref{hardware}, I propose several hardware implementations and discuss
design considerations, as well as the design process itself, that must be
considered for these designs. I also detail the results of some initial testing
investigating the performance expectations from differing hardware designs.

I conclude in \autoref{conclusions} by reflecting upon the results obtained
through this research, and making suggestions for further research in this
topic.

 
%%%%%%%%%%%%%%%%%%%%%%%%%%%%%%%%%%%%%%%%%%%%%%%%%
% Schedule
%%%%%%%%%%%%%%%%%%%%%%%%%%%%%%%%%%%%%%%%%%%%%%%%%
\section{Schedule}
\label{schedule}
A Gantt chart showing the anticipated schedule for the project is shown in 
\autoref{fig:ganttChart}. This will be updated as the project progresses.

% Gantt chart
% See http://www.martin-kumm.de/tex_gantt_package.php
\begin{figure}[h]
	\centering
	\scalebox{0.5}{
	\begin{gantt}[
			xunitlength=0.5cm,
			fontsize=\small,
			titlefontsize=\small,
			drawledgerline=true
			]{18}{36}
		% Months
		\begin{ganttitle}
			\titleelement{Mar}{4}
			\titleelement{Apr}{4}
			\titleelement{May}{4}
			\titleelement{Jun}{4}
			\titleelement{Jul}{4}
			\titleelement{Aug}{4}
			\titleelement{Sep}{4}
			\titleelement{Oct}{4}
			\titleelement{Nov}{4}
		\end{ganttitle}
		% Split months into quarters
		\begin{ganttitle}
			\numtitle{1}{1}{4}{1}
			\numtitle{1}{1}{4}{1}
			\numtitle{1}{1}{4}{1}
			\numtitle{1}{1}{4}{1}
			\numtitle{1}{1}{4}{1}
			\numtitle{1}{1}{4}{1}
			\numtitle{1}{1}{4}{1}
			\numtitle{1}{1}{4}{1}
			\numtitle{1}{1}{4}{1}
		\end{ganttitle}

		\input{ganttdata}
	\end{gantt}
	}
	\caption{Schedule for thesis work}
	\label{fig:ganttChart}
\end{figure}



%%%%%%%%%%%%%%%%%%%%%%%%%%%%%%%%%%%%%%%%%%%%%%%%%
% CHAPTER 04: Software Design and Evaluation
%%%%%%%%%%%%%%%%%%%%%%%%%%%%%%%%%%%%%%%%%%%%%%%%%
\chapter{Software Design and Evaluation}
\label{software}
\chapter{Introduction}
\label{introduction}

%%%%%%%%%%%%%%%%%%%%%%%%%%%%%%%%%%%%%%%%%%%%%%%%%
% Motivation
%%%%%%%%%%%%%%%%%%%%%%%%%%%%%%%%%%%%%%%%%%%%%%%%%
\section{Motivation}
\label{motivation}
Anomaly detection is an important technique which can be applied to a wide range
of applications. There are many techniques to detect and measure anomalies, each
usually most applicable to some specific problem domain \citeNeeded{}. A general
algorithm for anomaly detection has proved difficult to design \citeNeeded{},
due to various challenges associated with defining and measuring anomalies. Most
existing approaches are based on statistical or geometrical measures involving
distance. Whilst many algorithms work well for some subset of input data sets,
it has been difficult to discover an algorithm which performs both correctly and
efficiently on all possible data sets.

Previous research has found merit in applying randomization techniques to highly
multivariate data sets in order to reduce the dimensionality of these data sets
whilst maintaining their fundamental and statistical properties. Such reduction
of the dimensionality of data, assuming it can be performed efficiently, allows
previously-unscalable anomaly detection algorithms to be practically applied to
a wider range of data and applications.

Anomaly detection is an important and contemporary problem in the field of
computer science, and is of particular interesting to stock market analysis,
network intrusion detection and image comparison.

%%%%%%%%%%%%%%%%%%%%%%%%%%%%%%%%%%%%%%%%%%%%%%%%%
% Contributions of this thesis
%%%%%%%%%%%%%%%%%%%%%%%%%%%%%%%%%%%%%%%%%%%%%%%%%
\section{Contributions of this thesis}
\label{contributions}
One key difficulty in anomaly detection is the efficient scaling of a general
algorithm to apply to highly multivariate data \citeNeeded{}. In this thesis, I
explore the use of randomization techniques (such as random projections and
commute time) in anomaly detection algorithms. These techniques provide
encouraging results with regards to the run-time complexity of an algorithm.

Furthermore, in this thesis I attempt to make observations and an analysis of
the run-time performance of anomaly detection algorithms using randomization
techniques, so as to identify steps in the algorithms that bottleneck the
algorithm's performance. Through this identification it would be possible to
improve the \emph{actual} run-time performance of the algorithm by utilising
the advantages of reconfigurable computing.


%%%%%%%%%%%%%%%%%%%%%%%%%%%%%%%%%%%%%%%%%%%%%%%%%
% Organization
%%%%%%%%%%%%%%%%%%%%%%%%%%%%%%%%%%%%%%%%%%%%%%%%%
\section{Organization}
\label{organization}
The rest of this \thesis{} is organized as follows. In \autoref{background}, I
provide a background to various anomaly detection techniques, as well as
randomization techniques. In order to provide the reader with an understanding
of the background topics, a brief background is given to various topics in
linear algebra, vector calculus and graph theory. However, the provided
background material is not intended to be comprehensive, and it is expected that
the reader is familiar with these fields of mathematics. In
\autoref{reconfigurableComputing}, I provide an overview of reconfigurable
computing, including an explanation of \glspl{FPGA} and an overview of
\gls{HLS}.

In \autoref{software}, I profile the execution of an anomaly detection algorithm
and explore possible improvements to the algorithm, in particular by outsourcing
various stages of the algorithm to an FPGA device. I proceed to make iterative
refinements to the algorithm from its base implementation towards a hardware
design.

In \autoref{hardware}, I propose several hardware implementations and discuss
design considerations, as well as the design process itself, that must be
considered for these designs. I also detail the results of some initial testing
investigating the performance expectations from differing hardware designs.

I conclude in \autoref{conclusions} by reflecting upon the results obtained
through this research, and making suggestions for further research in this
topic.

 
%%%%%%%%%%%%%%%%%%%%%%%%%%%%%%%%%%%%%%%%%%%%%%%%%
% Schedule
%%%%%%%%%%%%%%%%%%%%%%%%%%%%%%%%%%%%%%%%%%%%%%%%%
\section{Schedule}
\label{schedule}
A Gantt chart showing the anticipated schedule for the project is shown in 
\autoref{fig:ganttChart}. This will be updated as the project progresses.

% Gantt chart
% See http://www.martin-kumm.de/tex_gantt_package.php
\begin{figure}[h]
	\centering
	\scalebox{0.5}{
	\begin{gantt}[
			xunitlength=0.5cm,
			fontsize=\small,
			titlefontsize=\small,
			drawledgerline=true
			]{18}{36}
		% Months
		\begin{ganttitle}
			\titleelement{Mar}{4}
			\titleelement{Apr}{4}
			\titleelement{May}{4}
			\titleelement{Jun}{4}
			\titleelement{Jul}{4}
			\titleelement{Aug}{4}
			\titleelement{Sep}{4}
			\titleelement{Oct}{4}
			\titleelement{Nov}{4}
		\end{ganttitle}
		% Split months into quarters
		\begin{ganttitle}
			\numtitle{1}{1}{4}{1}
			\numtitle{1}{1}{4}{1}
			\numtitle{1}{1}{4}{1}
			\numtitle{1}{1}{4}{1}
			\numtitle{1}{1}{4}{1}
			\numtitle{1}{1}{4}{1}
			\numtitle{1}{1}{4}{1}
			\numtitle{1}{1}{4}{1}
			\numtitle{1}{1}{4}{1}
		\end{ganttitle}

		\input{ganttdata}
	\end{gantt}
	}
	\caption{Schedule for thesis work}
	\label{fig:ganttChart}
\end{figure}



%%%%%%%%%%%%%%%%%%%%%%%%%%%%%%%%%%%%%%%%%%%%%%%%%
% CHAPTER 05: Hardware Design and Implementation
%%%%%%%%%%%%%%%%%%%%%%%%%%%%%%%%%%%%%%%%%%%%%%%%%
\chapter{Hardware Design and Implementation}
\label{hardware}
\chapter{Introduction}
\label{introduction}

%%%%%%%%%%%%%%%%%%%%%%%%%%%%%%%%%%%%%%%%%%%%%%%%%
% Motivation
%%%%%%%%%%%%%%%%%%%%%%%%%%%%%%%%%%%%%%%%%%%%%%%%%
\section{Motivation}
\label{motivation}
Anomaly detection is an important technique which can be applied to a wide range
of applications. There are many techniques to detect and measure anomalies, each
usually most applicable to some specific problem domain \citeNeeded{}. A general
algorithm for anomaly detection has proved difficult to design \citeNeeded{},
due to various challenges associated with defining and measuring anomalies. Most
existing approaches are based on statistical or geometrical measures involving
distance. Whilst many algorithms work well for some subset of input data sets,
it has been difficult to discover an algorithm which performs both correctly and
efficiently on all possible data sets.

Previous research has found merit in applying randomization techniques to highly
multivariate data sets in order to reduce the dimensionality of these data sets
whilst maintaining their fundamental and statistical properties. Such reduction
of the dimensionality of data, assuming it can be performed efficiently, allows
previously-unscalable anomaly detection algorithms to be practically applied to
a wider range of data and applications.

Anomaly detection is an important and contemporary problem in the field of
computer science, and is of particular interesting to stock market analysis,
network intrusion detection and image comparison.

%%%%%%%%%%%%%%%%%%%%%%%%%%%%%%%%%%%%%%%%%%%%%%%%%
% Contributions of this thesis
%%%%%%%%%%%%%%%%%%%%%%%%%%%%%%%%%%%%%%%%%%%%%%%%%
\section{Contributions of this thesis}
\label{contributions}
One key difficulty in anomaly detection is the efficient scaling of a general
algorithm to apply to highly multivariate data \citeNeeded{}. In this thesis, I
explore the use of randomization techniques (such as random projections and
commute time) in anomaly detection algorithms. These techniques provide
encouraging results with regards to the run-time complexity of an algorithm.

Furthermore, in this thesis I attempt to make observations and an analysis of
the run-time performance of anomaly detection algorithms using randomization
techniques, so as to identify steps in the algorithms that bottleneck the
algorithm's performance. Through this identification it would be possible to
improve the \emph{actual} run-time performance of the algorithm by utilising
the advantages of reconfigurable computing.


%%%%%%%%%%%%%%%%%%%%%%%%%%%%%%%%%%%%%%%%%%%%%%%%%
% Organization
%%%%%%%%%%%%%%%%%%%%%%%%%%%%%%%%%%%%%%%%%%%%%%%%%
\section{Organization}
\label{organization}
The rest of this \thesis{} is organized as follows. In \autoref{background}, I
provide a background to various anomaly detection techniques, as well as
randomization techniques. In order to provide the reader with an understanding
of the background topics, a brief background is given to various topics in
linear algebra, vector calculus and graph theory. However, the provided
background material is not intended to be comprehensive, and it is expected that
the reader is familiar with these fields of mathematics. In
\autoref{reconfigurableComputing}, I provide an overview of reconfigurable
computing, including an explanation of \glspl{FPGA} and an overview of
\gls{HLS}.

In \autoref{software}, I profile the execution of an anomaly detection algorithm
and explore possible improvements to the algorithm, in particular by outsourcing
various stages of the algorithm to an FPGA device. I proceed to make iterative
refinements to the algorithm from its base implementation towards a hardware
design.

In \autoref{hardware}, I propose several hardware implementations and discuss
design considerations, as well as the design process itself, that must be
considered for these designs. I also detail the results of some initial testing
investigating the performance expectations from differing hardware designs.

I conclude in \autoref{conclusions} by reflecting upon the results obtained
through this research, and making suggestions for further research in this
topic.

 
%%%%%%%%%%%%%%%%%%%%%%%%%%%%%%%%%%%%%%%%%%%%%%%%%
% Schedule
%%%%%%%%%%%%%%%%%%%%%%%%%%%%%%%%%%%%%%%%%%%%%%%%%
\section{Schedule}
\label{schedule}
A Gantt chart showing the anticipated schedule for the project is shown in 
\autoref{fig:ganttChart}. This will be updated as the project progresses.

% Gantt chart
% See http://www.martin-kumm.de/tex_gantt_package.php
\begin{figure}[h]
	\centering
	\scalebox{0.5}{
	\begin{gantt}[
			xunitlength=0.5cm,
			fontsize=\small,
			titlefontsize=\small,
			drawledgerline=true
			]{18}{36}
		% Months
		\begin{ganttitle}
			\titleelement{Mar}{4}
			\titleelement{Apr}{4}
			\titleelement{May}{4}
			\titleelement{Jun}{4}
			\titleelement{Jul}{4}
			\titleelement{Aug}{4}
			\titleelement{Sep}{4}
			\titleelement{Oct}{4}
			\titleelement{Nov}{4}
		\end{ganttitle}
		% Split months into quarters
		\begin{ganttitle}
			\numtitle{1}{1}{4}{1}
			\numtitle{1}{1}{4}{1}
			\numtitle{1}{1}{4}{1}
			\numtitle{1}{1}{4}{1}
			\numtitle{1}{1}{4}{1}
			\numtitle{1}{1}{4}{1}
			\numtitle{1}{1}{4}{1}
			\numtitle{1}{1}{4}{1}
			\numtitle{1}{1}{4}{1}
		\end{ganttitle}

		\input{ganttdata}
	\end{gantt}
	}
	\caption{Schedule for thesis work}
	\label{fig:ganttChart}
\end{figure}



%%%%%%%%%%%%%%%%%%%%%%%%%%%%%%%%%%%%%%%%%%%%%%%%%
% CHAPTER 06: Conclusions and Future Work
%%%%%%%%%%%%%%%%%%%%%%%%%%%%%%%%%%%%%%%%%%%%%%%%%
\chapter{Conclusions and Future Work}
\label{conclusions}
\chapter{Introduction}
\label{introduction}

%%%%%%%%%%%%%%%%%%%%%%%%%%%%%%%%%%%%%%%%%%%%%%%%%
% Motivation
%%%%%%%%%%%%%%%%%%%%%%%%%%%%%%%%%%%%%%%%%%%%%%%%%
\section{Motivation}
\label{motivation}
Anomaly detection is an important technique which can be applied to a wide range
of applications. There are many techniques to detect and measure anomalies, each
usually most applicable to some specific problem domain \citeNeeded{}. A general
algorithm for anomaly detection has proved difficult to design \citeNeeded{},
due to various challenges associated with defining and measuring anomalies. Most
existing approaches are based on statistical or geometrical measures involving
distance. Whilst many algorithms work well for some subset of input data sets,
it has been difficult to discover an algorithm which performs both correctly and
efficiently on all possible data sets.

Previous research has found merit in applying randomization techniques to highly
multivariate data sets in order to reduce the dimensionality of these data sets
whilst maintaining their fundamental and statistical properties. Such reduction
of the dimensionality of data, assuming it can be performed efficiently, allows
previously-unscalable anomaly detection algorithms to be practically applied to
a wider range of data and applications.

Anomaly detection is an important and contemporary problem in the field of
computer science, and is of particular interesting to stock market analysis,
network intrusion detection and image comparison.

%%%%%%%%%%%%%%%%%%%%%%%%%%%%%%%%%%%%%%%%%%%%%%%%%
% Contributions of this thesis
%%%%%%%%%%%%%%%%%%%%%%%%%%%%%%%%%%%%%%%%%%%%%%%%%
\section{Contributions of this thesis}
\label{contributions}
One key difficulty in anomaly detection is the efficient scaling of a general
algorithm to apply to highly multivariate data \citeNeeded{}. In this thesis, I
explore the use of randomization techniques (such as random projections and
commute time) in anomaly detection algorithms. These techniques provide
encouraging results with regards to the run-time complexity of an algorithm.

Furthermore, in this thesis I attempt to make observations and an analysis of
the run-time performance of anomaly detection algorithms using randomization
techniques, so as to identify steps in the algorithms that bottleneck the
algorithm's performance. Through this identification it would be possible to
improve the \emph{actual} run-time performance of the algorithm by utilising
the advantages of reconfigurable computing.


%%%%%%%%%%%%%%%%%%%%%%%%%%%%%%%%%%%%%%%%%%%%%%%%%
% Organization
%%%%%%%%%%%%%%%%%%%%%%%%%%%%%%%%%%%%%%%%%%%%%%%%%
\section{Organization}
\label{organization}
The rest of this \thesis{} is organized as follows. In \autoref{background}, I
provide a background to various anomaly detection techniques, as well as
randomization techniques. In order to provide the reader with an understanding
of the background topics, a brief background is given to various topics in
linear algebra, vector calculus and graph theory. However, the provided
background material is not intended to be comprehensive, and it is expected that
the reader is familiar with these fields of mathematics. In
\autoref{reconfigurableComputing}, I provide an overview of reconfigurable
computing, including an explanation of \glspl{FPGA} and an overview of
\gls{HLS}.

In \autoref{software}, I profile the execution of an anomaly detection algorithm
and explore possible improvements to the algorithm, in particular by outsourcing
various stages of the algorithm to an FPGA device. I proceed to make iterative
refinements to the algorithm from its base implementation towards a hardware
design.

In \autoref{hardware}, I propose several hardware implementations and discuss
design considerations, as well as the design process itself, that must be
considered for these designs. I also detail the results of some initial testing
investigating the performance expectations from differing hardware designs.

I conclude in \autoref{conclusions} by reflecting upon the results obtained
through this research, and making suggestions for further research in this
topic.

 
%%%%%%%%%%%%%%%%%%%%%%%%%%%%%%%%%%%%%%%%%%%%%%%%%
% Schedule
%%%%%%%%%%%%%%%%%%%%%%%%%%%%%%%%%%%%%%%%%%%%%%%%%
\section{Schedule}
\label{schedule}
A Gantt chart showing the anticipated schedule for the project is shown in 
\autoref{fig:ganttChart}. This will be updated as the project progresses.

% Gantt chart
% See http://www.martin-kumm.de/tex_gantt_package.php
\begin{figure}[h]
	\centering
	\scalebox{0.5}{
	\begin{gantt}[
			xunitlength=0.5cm,
			fontsize=\small,
			titlefontsize=\small,
			drawledgerline=true
			]{18}{36}
		% Months
		\begin{ganttitle}
			\titleelement{Mar}{4}
			\titleelement{Apr}{4}
			\titleelement{May}{4}
			\titleelement{Jun}{4}
			\titleelement{Jul}{4}
			\titleelement{Aug}{4}
			\titleelement{Sep}{4}
			\titleelement{Oct}{4}
			\titleelement{Nov}{4}
		\end{ganttitle}
		% Split months into quarters
		\begin{ganttitle}
			\numtitle{1}{1}{4}{1}
			\numtitle{1}{1}{4}{1}
			\numtitle{1}{1}{4}{1}
			\numtitle{1}{1}{4}{1}
			\numtitle{1}{1}{4}{1}
			\numtitle{1}{1}{4}{1}
			\numtitle{1}{1}{4}{1}
			\numtitle{1}{1}{4}{1}
			\numtitle{1}{1}{4}{1}
		\end{ganttitle}

		\input{ganttdata}
	\end{gantt}
	}
	\caption{Schedule for thesis work}
	\label{fig:ganttChart}
\end{figure}



%%%%%%%%%%%%%%%%%%%%%%%%%%%%%%%%%%%%%%%%%%%%%%%%%
% Appendix
%%%%%%%%%%%%%%%%%%%%%%%%%%%%%%%%%%%%%%%%%%%%%%%%%
\appendix
\chapter{Introduction}
\label{introduction}

%%%%%%%%%%%%%%%%%%%%%%%%%%%%%%%%%%%%%%%%%%%%%%%%%
% Motivation
%%%%%%%%%%%%%%%%%%%%%%%%%%%%%%%%%%%%%%%%%%%%%%%%%
\section{Motivation}
\label{motivation}
Anomaly detection is an important technique which can be applied to a wide range
of applications. There are many techniques to detect and measure anomalies, each
usually most applicable to some specific problem domain \citeNeeded{}. A general
algorithm for anomaly detection has proved difficult to design \citeNeeded{},
due to various challenges associated with defining and measuring anomalies. Most
existing approaches are based on statistical or geometrical measures involving
distance. Whilst many algorithms work well for some subset of input data sets,
it has been difficult to discover an algorithm which performs both correctly and
efficiently on all possible data sets.

Previous research has found merit in applying randomization techniques to highly
multivariate data sets in order to reduce the dimensionality of these data sets
whilst maintaining their fundamental and statistical properties. Such reduction
of the dimensionality of data, assuming it can be performed efficiently, allows
previously-unscalable anomaly detection algorithms to be practically applied to
a wider range of data and applications.

Anomaly detection is an important and contemporary problem in the field of
computer science, and is of particular interesting to stock market analysis,
network intrusion detection and image comparison.

%%%%%%%%%%%%%%%%%%%%%%%%%%%%%%%%%%%%%%%%%%%%%%%%%
% Contributions of this thesis
%%%%%%%%%%%%%%%%%%%%%%%%%%%%%%%%%%%%%%%%%%%%%%%%%
\section{Contributions of this thesis}
\label{contributions}
One key difficulty in anomaly detection is the efficient scaling of a general
algorithm to apply to highly multivariate data \citeNeeded{}. In this thesis, I
explore the use of randomization techniques (such as random projections and
commute time) in anomaly detection algorithms. These techniques provide
encouraging results with regards to the run-time complexity of an algorithm.

Furthermore, in this thesis I attempt to make observations and an analysis of
the run-time performance of anomaly detection algorithms using randomization
techniques, so as to identify steps in the algorithms that bottleneck the
algorithm's performance. Through this identification it would be possible to
improve the \emph{actual} run-time performance of the algorithm by utilising
the advantages of reconfigurable computing.


%%%%%%%%%%%%%%%%%%%%%%%%%%%%%%%%%%%%%%%%%%%%%%%%%
% Organization
%%%%%%%%%%%%%%%%%%%%%%%%%%%%%%%%%%%%%%%%%%%%%%%%%
\section{Organization}
\label{organization}
The rest of this \thesis{} is organized as follows. In \autoref{background}, I
provide a background to various anomaly detection techniques, as well as
randomization techniques. In order to provide the reader with an understanding
of the background topics, a brief background is given to various topics in
linear algebra, vector calculus and graph theory. However, the provided
background material is not intended to be comprehensive, and it is expected that
the reader is familiar with these fields of mathematics. In
\autoref{reconfigurableComputing}, I provide an overview of reconfigurable
computing, including an explanation of \glspl{FPGA} and an overview of
\gls{HLS}.

In \autoref{software}, I profile the execution of an anomaly detection algorithm
and explore possible improvements to the algorithm, in particular by outsourcing
various stages of the algorithm to an FPGA device. I proceed to make iterative
refinements to the algorithm from its base implementation towards a hardware
design.

In \autoref{hardware}, I propose several hardware implementations and discuss
design considerations, as well as the design process itself, that must be
considered for these designs. I also detail the results of some initial testing
investigating the performance expectations from differing hardware designs.

I conclude in \autoref{conclusions} by reflecting upon the results obtained
through this research, and making suggestions for further research in this
topic.

 
%%%%%%%%%%%%%%%%%%%%%%%%%%%%%%%%%%%%%%%%%%%%%%%%%
% Schedule
%%%%%%%%%%%%%%%%%%%%%%%%%%%%%%%%%%%%%%%%%%%%%%%%%
\section{Schedule}
\label{schedule}
A Gantt chart showing the anticipated schedule for the project is shown in 
\autoref{fig:ganttChart}. This will be updated as the project progresses.

% Gantt chart
% See http://www.martin-kumm.de/tex_gantt_package.php
\begin{figure}[h]
	\centering
	\scalebox{0.5}{
	\begin{gantt}[
			xunitlength=0.5cm,
			fontsize=\small,
			titlefontsize=\small,
			drawledgerline=true
			]{18}{36}
		% Months
		\begin{ganttitle}
			\titleelement{Mar}{4}
			\titleelement{Apr}{4}
			\titleelement{May}{4}
			\titleelement{Jun}{4}
			\titleelement{Jul}{4}
			\titleelement{Aug}{4}
			\titleelement{Sep}{4}
			\titleelement{Oct}{4}
			\titleelement{Nov}{4}
		\end{ganttitle}
		% Split months into quarters
		\begin{ganttitle}
			\numtitle{1}{1}{4}{1}
			\numtitle{1}{1}{4}{1}
			\numtitle{1}{1}{4}{1}
			\numtitle{1}{1}{4}{1}
			\numtitle{1}{1}{4}{1}
			\numtitle{1}{1}{4}{1}
			\numtitle{1}{1}{4}{1}
			\numtitle{1}{1}{4}{1}
			\numtitle{1}{1}{4}{1}
		\end{ganttitle}

		\input{ganttdata}
	\end{gantt}
	}
	\caption{Schedule for thesis work}
	\label{fig:ganttChart}
\end{figure}



%%%%%%%%%%%%%%%%%%%%%%%%%%%%%%%%%%%%%%%%%%%%%%%%%
% Glossary
%%%%%%%%%%%%%%%%%%%%%%%%%%%%%%%%%%%%%%%%%%%%%%%%%
\glsaddall
\printglossary

%%%%%%%%%%%%%%%%%%%%%%%%%%%%%%%%%%%%%%%%%%%%%%%%%
% Bibliography
%%%%%%%%%%%%%%%%%%%%%%%%%%%%%%%%%%%%%%%%%%%%%%%%%
\nocite{*}
\printbibliography

\end{document}
